\section{Uniqueness of the five incidence patterns}\label{sec:maps}

The graph of a convex polyhedron is simple: there are no loops or
multiple edges. Two distinct facets share at most one edge, since
their intersection is a convex subset of their common line. We use
the cyclic order at each vertex, supplied by its convex vertex figure.

For $(p,q)=(3,3)$, there are four vertices of degree three. The graph
is $K_4$, and its four triples are precisely the four faces. This is
the tetrahedral pattern.

For $(p,q)=(3,4)$, there are six vertices of degree four. Each vertex
has exactly one nonneighbor among the other five vertices, so
nonadjacency pairs the vertices into three disjoint pairs. The graph
is $K_{2,2,2}$. Its triangular cliques choose one vertex from each
pair, giving eight triples. There are eight faces, so all these
triples are faces. This is the octahedral pattern.

\subsection{The twelve-vertex, valence-five case}

Here all faces are triangles, every vertex has degree five, and
$V=12$. We first show that \emph{every three-cycle is a face}.
Otherwise choose a nonfacial three-cycle. By the triangular separation
argument of Section~\ref{sec:euler}, it has $k>0$ vertices on one side
and $l>0$ on the other, with $k+l=9$. No edge crosses its boundary,
so interior vertices retain their full degree five. Applying
Lemma~\ref{lem:disk} on either side gives $e_{II}=2k$. Simplicity gives
\[
 2k\leq\binom k2.
\]
The four cases $k=1,2,3,4$ violate this inequality; hence $k\geq5$.
Similarly $l\geq5$, contradicting $k+l=9$.

Fix a vertex $A$, with cyclic neighbors $v_0,\ldots,v_4$, subscripts
read modulo five. Consecutive neighbors are joined. Nonconsecutive
neighbors cannot be joined, since that would form a nonfacial
three-cycle with $A$. Thus the neighbors induce exactly a five-cycle.
Across the edge $v_iv_{i+1}$ from the face $Av_iv_{i+1}$, let the other
face have third vertex $w_i$. Each $w_i$ lies outside
$\{A,v_0,\ldots,v_4\}$: otherwise that face supplies a forbidden
chord of the five-cycle.

Adjacent $w_i$ are distinct. If $w_i=w_{i+1}=w$, the link of
$v_{i+1}$ is already the closed four-cycle
$A,v_i,w,v_{i+2}$, contradicting its degree five. Nonadjacent $w_i$
are also distinct. After reversing the indexing if necessary, a
nonadjacent coincidence is $w_i=w_{i+2}=w$. Then
$wv_{i+1}v_{i+2}$ is a three-cycle. It is not a face, because the two
faces on $v_{i+1}v_{i+2}$ have third vertices $A$ and $w_{i+1}$,
both different from $w$. This is again impossible.

We now have eleven distinct vertices $A,v_0,\ldots,v_4,w_0,\ldots,w_4$;
call the remaining vertex $B$. At $v_i$, the five already known
neighbors are
\[
 A,\ v_{i-1},\ v_{i+1},\ w_{i-1},\ w_i.
\]
The cyclic link is exhausted, forcing the face $v_iw_{i-1}w_i$ and
the edge $w_{i-1}w_i$. Neither $A$ nor any $v_i$ can have another
neighbor. Therefore the five neighbors of $B$ must be the five
$w_i$. Each $w_i$ has now acquired precisely five neighbors, and
the full list of twenty faces is
\[
 Av_iv_{i+1},\quad v_iv_{i+1}w_i,\quad
 v_iw_{i-1}w_i,\quad Bw_{i-1}w_i\qquad(0\leq i<5).
\]
This determines the icosahedral pattern, including its face cycles.

\subsection{Duality supplies the other two cases}

Interchange vertices and facets, keeping an edge for each old edge.
The new cyclic incidences are the old facet boundaries and vertex
links. This finite dual map is simple and is again a spherical
polyhedral map. One can draw it explicitly by joining each facet
center to its edge midpoints on the old surface; the two halves
through each old edge form a new edge. The local polygons fill the
vertex stars, with no crossing. Equivalently, translate an interior
point to zero and write the facet inequalities as $a_f\cdot x\leq1$;
the polar polyhedron has vertices $a_f$ and these reversed incidences.
Here is a finite verification of the latter assertion. Define $Q$ by
$v\cdot y\leq1$ for every old vertex $v$. An explicit small cube around
zero lies in the original polyhedron, so these inequalities bound all
coordinates of $Q$; conversely a sufficiently small cube lies in $Q$.
A vertex $y$ of $Q$ has three linearly independent tight old vertices:
otherwise a nonzero null direction and the finite list of strict slacks
give two opposite feasible perturbations, contradicting extremality.
Those vertices are affinely independent on the support plane
$y\cdot x=1$, so that plane is an old facet and $y=a_f$.
Conversely three affinely independent vertices of an old facet make
$a_f$ a vertex of $Q$. At an old vertex $v$, its incident facet normals
span three dimensions by the same perturbation test; their normalized
points span the facet $v\cdot y=1$ of $Q$. Tight-incidence equations
therefore reverse all incidences, including edges. This proves duality
using finite rank and slack computations.

The dual of the $(4,3)$ case is the already classified $(3,4)$ case,
so its map is the cube. The dual of the $(5,3)$ case is the already
classified $(3,5)$ case, so its map is the dodecahedron. Notice that
we only use \emph{combinatorial} duality here; we do not assume in
advance that the metric dual has regular faces.

\section{Finite geometric rigidity}\label{sec:rigidity}

Incidence uniqueness alone does not establish uniqueness of shape.
We now prove the needed version of Cauchy's rigidity theorem. The
argument is expressed using exact scalar comparisons. The corrected
arm induction is the one of Schoenberg--Zaremba~\cite{sz}; we supply
the finite arithmetic interpretation of its geometric constructions.
The overall sign argument is the classical Cauchy method, also
reviewed in~\cite{bowers}.

\subsection{Elementary spherical comparisons}

At a vertex, normalize the outgoing edge vectors to unit vectors.
They lie in an open hemisphere and bound a convex spherical polygon.
The sides encode the face angles and its angles encode the dihedral
angles. These are finite vectors obtained by adjoining positive
square roots. An angle or spherical distance is represented by its
cosine, together with the sign of its positive sine when needed.
On the interval from zero to a straight angle, comparison of angles
is the reversed comparison of their cosines. ``Angle'' notation in
this proof abbreviates these ordered unit-circle data, not an appeal
to an analytic inverse cosine function.

Resolving two unit vectors into their components parallel and
perpendicular to a third gives the spherical hinge identity
\begin{equation}\label{eq:hinge}
 \cos c=\cos a\cos b+\sin a\sin b\cos\gamma.
\end{equation}
Here $a,b\in(0,\pi)$ are the fixed incident short sides, $c$ the
opposite short side, and $\gamma\in(0,\pi)$ the included angle.
Since $\sin a\sin b>0$, increasing $\gamma$ strictly increases $c$.
This is an ordered-field dot-product calculation.

The short-arc triangle inequality has an equally algebraic proof.
For unit $u,v,w$, put $a=d(u,v)$, $b=d(v,w)$. Splitting $u,w$ into
their components along $v$ yields
\[
 u\cdot w\geq\cos a\cos b-\sin a\sin b=\cos(a+b).
\]
If $a+b\leq\pi$, the reversed cosine order gives
$d(u,w)\leq a+b$; if $a+b\geq\pi$, this is immediate from the
definition of short distance. The required dot-product inequality
follows from the nonnegativity of a squared norm. Angle addition is
the multiplication law for pairs $(\cos a,\sin a)$, with finite
sign cases selecting the half-circle. Thus these facts do not need
real completeness or calculus.

\begin{lemma}[Spherical arm comparison]\label{lem:arm}
Let $P_1,\ldots,P_n$ and $Q_1,\ldots,Q_n$ be strictly convex spherical
polygons in open hemispheres. Suppose the successive sides for
$1\leq i<n$ have equal corresponding lengths, and every internal
chain angle at $P_i$, $2\leq i<n$, is at most that at $Q_i$.
Then $d(P_1,P_n)\leq d(Q_1,Q_n)$, strictly if any internal angle
increases.
\end{lemma}
\begin{proof}
The case $n=3$ is~\eqref{eq:hinge}. Induct on $n$. If an internal
angle agrees, cut off the triangle at that vertex in both polygons.
The removed triangles are congruent by two sides and the included
angle. The new corresponding diagonal lengths agree, and identical
triangle angles are subtracted from the neighboring polygon angles.
All remaining inequalities, including any strict one, are preserved.
Apply induction to the remaining convex polygons.

Otherwise all internal angles increase strictly. Keep
$P_1,\ldots,P_{n-1}$ fixed and open only the angle at $P_{n-1}$,
preserving its distance to the moving last endpoint. Stop at its
target angle or the first flat corner. The following paragraph
specifies this stop by a finite calculation. Before it, the angle
at $P_{n-1}$ in the triangle $P_1P_{n-1}P_n$ increases by the same
amount, since the fixed diagonal lies inside the polygon wedge.
The hinge identity therefore strictly increases the endpoint distance.

Only the corners at $P_1$, $P_{n-1}$, and the endpoint can change.
The angle at $P_{n-1}$ remains below $\pi$, as does its target.
The endpoint cannot flatten: it would lie between $P_1$ and
$P_{n-1}$ on their short arc, making the angle at $P_{n-1}$ in that
triangle zero, although it has increased from a positive value.
Thus an earlier flat corner must be $P_1$.

If the target is reached without a flat corner, cut off the now
congruent last triangles and apply the equal-angle reduction. If a
flat corner is reached, write $X$ for the endpoint. Its great-circle
order is $P_2,P_1,X$, so
\[
 d(P_2,X)=d(P_2,P_1)+d(P_1,X).
\]
Remove $P_1$ and apply induction to the smaller convex polygons
$P_2,\ldots,P_{n-1},X$ and $Q_2,\ldots,Q_n$. Their chain sides agree
and their internal angles have the required inequalities. By the
triangle inequality,
\begin{align*}
 d(Q_1,Q_n)&\geq d(Q_2,Q_n)-d(Q_1,Q_2)\\
 &\geq d(P_2,X)-d(P_1,P_2)\\
 &=d(P_1,X)>d(P_1,P_n).
\end{align*}
The last strict inequality was obtained before the flat-corner
reduction. It also handles a target and flat-corner event occurring
simultaneously. Flat auxiliary corners are removed immediately, so
induction is applied to strictly convex polygons.
\end{proof}

\paragraph{How to compute the first event.}
Let $u=P_{n-1}$ and $x_0=P_n$ be unit vectors. Put
\[
 c=u\cdot x_0,\quad s=\sqrt{1-c^2}>0,\quad
 e=(x_0-cu)/s,\quad f=\pm(u\times e),
\]
choosing the sign of $f$ that opens the angle. The relevant arc of the
endpoint circle is parametrized by
\begin{equation}\label{eq:circleparam}
 x(t)=cu+s\frac{(1-t^2)e+2tf}{1+t^2}.
\end{equation}
For initial and target pivot angles $\alpha_0<\alpha_1<\pi$, put
\[
 C=\cos\alpha_1\cos\alpha_0+\sin\alpha_1\sin\alpha_0,
 \qquad
 S=\sin\alpha_1\cos\alpha_0-\cos\alpha_1\sin\alpha_0.
\]
The target is $t_*=S/(1+C)>0$, a finite constructible scalar.
Every supporting-side determinant is
a polynomial of degree at most two in $t$, divided by the positive
$1+t^2$. Solve the finitely many linear or quadratic equations,
retain roots on the directed target interval, and choose the first,
or the target if it is earlier. Positivity before this event follows
from factoring these polynomials and testing their signs; no infimum
or intermediate value theorem is being invoked.

For completeness, the boundary check used in the proof is finite as
well. At a convex polygon, the vertices on one supporting great
circle form a consecutive boundary interval. A first support equality
at a nonconsecutive vertex would therefore flatten an intermediate
corner. All such fixed internal corners are strictly below $\pi$.
Coincidence of the two endpoints is excluded by the strictly increasing
positive endpoint distance. An open hemisphere persists even at the
first event. To verify this without a limiting argument, let $h$ be
the sum of the inward oriented side normals at that event. All their
vertex dot products are nonnegative. For $n\geq5$, every fixed vertex
has a nonincident fixed side with strictly positive support. Moreover
the three fixed normals $P_1\times P_2$, $P_2\times P_3$,
$P_3\times P_4$ span three dimensions: the first two span
$P_2^\perp$, while the third has positive dot product with $P_2$.
The moving endpoint therefore cannot have zero dot product with all
the fixed normals. Hence $h$ has positive dot product with every
vertex. For $n=4$, the two fixed normals have common kernel
$\operatorname{span}(P_2)$. Failure of positivity at $P_2$ or at the
endpoint would force the endpoint to be $P_2$ or $-P_2$, making the
opened pivot angle zero or $\pi$, both excluded. The other fixed
vertices already have a strictly positive fixed-side support.

Projection $x\mapsto x/(h\cdot x)$ thus identifies the event polygon
with a weakly convex ordinary polygon in the plane $h\cdot x=1$.
An extra support equality makes a consecutive boundary chain flat.
A repeated non-end vertex would collapse such a chain and also
flatten a fixed internal corner. Thus one of the three changing
corners must flatten; the preceding exclusions leave only $P_1$.
This finite boundary analysis justifies the alternative used in the
arm proof and prevents a loss of the hemisphere condition from being
hidden in the geometric wording.

\begin{lemma}[Four sign changes]\label{lem:four}
For two convex spherical polygons with equal corresponding sides,
label each vertex by the sign of the difference of corresponding
angles. Either every sign is zero, or the cyclic sequence obtained
by deleting zeros has at least four sign changes.
\end{lemma}
\begin{proof}
A cyclic sign sequence has an even number of changes. If all nonzero
signs agree, cut the polygon at an edge not incident to a chosen
strictly changed vertex. Lemma~\ref{lem:arm} makes the distance
between its endpoints change strictly, although that edge length
is fixed. If there are exactly two changes, choose corresponding
midpoints on edges in the two transition gaps. The diagonal between
them splits each polygon into two convex chains. On one chain the
angle comparisons all have one sign, with a strict comparison; on
the other they have the opposite sign. The arm lemma makes the
same diagonal both longer and shorter. Both alternatives are
impossible. Midpoints are normalized sums of unit vectors, so this
construction also uses only square roots and field operations.
\end{proof}

\subsection{A sign-count argument using finite permutations}

We give the combinatorial ingredient without a hidden planar
separation theorem. A \emph{rotation graph} consists of its directed
edges, the reversal involution, and the permutation giving the cyclic
order at each vertex. The cycles obtained by turning at each vertex
and traversing the next edge are the face walks. An isolated vertex
is assigned one empty face walk. All these data are finite lists and
permutations. For the polyhedral rotation graph they are the actual
facet boundaries, and Euler's identity has already been proved.

\begin{lemma}[Euler inheritance under edge deletion]\label{lem:rotation}
Every connected finite rotation graph has $V-E+F\leq2$. If equality
holds, deleting any collection of edges leaves each nonempty connected
component, with its induced cyclic orders, satisfying $V-E+F=2$.
\end{lemma}
\begin{proof}
Contract a spanning tree. A nonloop contraction merges its two vertex
cycles and preserves face walks, while reducing both $V$ and $E$ by
one. The result is one vertex with $m=E-V+1$ loops. Start with one
vertex and its empty face, and insert these loops successively.
Each insertion either merges two face cycles or splits one, so the
final face count is at most $m+1$. Therefore $V-E+F\leq2$.

Now start with equality and delete one edge. If it is not a bridge,
the result is connected. The face-cycle splice changes $F$ by $-1$
or $+1$, making its Euler value two or four. The bound just proved
excludes four. If the edge is a bridge, its two directions lie on
the same face walk: after entering the separated component that walk
must eventually leave across the reverse direction. Deletion splits
this walk into one in each component, with the empty-face convention
covering isolated vertices. The sum of the two Euler values is four.
Each is at most two, so each is two. Iterate these finite operations;
discarding isolated components does not affect the others.
\end{proof}

\begin{lemma}[No everywhere-four-change labeling]\label{lem:signcount}
A nonempty subgraph of a simple polyhedral graph cannot have signed
edges with at least four cyclic sign changes at every vertex incident
to an edge.
\end{lemma}
\begin{proof}
Choose a connected component of the signed graph and delete everything
else. By Lemma~\ref{lem:rotation}, $V-E+F=2$ for its face walks.
Every vertex has degree at least four. Each face walk has length
$d\geq3$: length one requires a loop, and length two in a simple
connected graph requires the single-edge component, whose endpoints
have degree one. A face walk of length three has at most two sign
changes. For $d\geq4$, it has at most $d\leq2d-4$. Count changes
at the corners, either by vertices or by face walks. Since the total
face-walk length is $2E$, the resulting count $S$ satisfies
\[
 4V\leq S\leq\sum_f(2d_f-4)=4E-4F=4V-8,
\]
a contradiction. Bridges and repeated vertices in face walks cause
no difficulty: the corner count is on directed-edge occurrences.
\end{proof}

\begin{theorem}[Finite Cauchy rigidity]\label{thm:cauchy}
Two full-dimensional convex polyhedra with isomorphic face incidences
and congruent corresponding facets are congruent.
\end{theorem}
\begin{proof}
For each corresponding edge, compare its two interior dihedral
angles and assign $+,0,-$. This comparison is decidable: the cosine
of an interior dihedral angle is
$-n_1\cdot n_2/(\lVert n_1\rVert\lVert n_2\rVert)$ for outward
facet normals. At corresponding vertices the spherical links have
equal side lengths, namely the corresponding facet angles.
Lemma~\ref{lem:four} therefore says that any vertex incident to a
nonzero sign has at least four cyclic sign changes after zeros are
omitted. Delete the zero edges. Lemma~\ref{lem:signcount} forces the
remaining graph to be empty. All corresponding dihedral angles agree.

Align one pair of congruent facets and place both polyhedron interiors
on the same side of their plane. A neighboring facet is determined
by its common edge, its dihedral angle, and the side prescribed by
convexity. It must align too. Traverse a finite spanning tree of the
face adjacency graph. Every facet and hence every vertex aligns.
The resulting Euclidean isometry is the required congruence.
\end{proof}

\paragraph{Uniqueness of the small regular face metrics.}
One may verify the remaining elementary polygon fact without
trigonometric functions. Center a regular face and let $R$ be its
cyclic rotation. In oriented orthonormal coordinates,
$R=\left(\begin{smallmatrix}c&-s\\s&c\end{smallmatrix}\right)$,
$c^2+s^2=1$, $s>0$, and its order is exactly $p$. Put $t=2c$;
then $R^2-tR+I=0$. For $p=3$, the identity
$I+R+R^2=0$ gives $t=-1$. For $p=4$, exact order four gives
$R^2=-I$, hence $t=0$. These calculations take place in the
two-dimensional algebra of rotation matrices, where a nonzero
element $aI+bJ$, $J^2=-I$, is invertible by its positive norm.
For $p=5$, divide $R^5-I$ by the invertible $R-I$ and multiply
by $R^{-2}$. Since $R+R^{-1}=tI$, the result is
\[
 t^2+t-1=0.
\]
The negative root is less than $-1$ and violates convex order:
for unit centered vertices $u_i=R^iu_0$,
\[
 \det(u_1-u_0,u_3-u_0)=2s(1-c)(t+1)>0
\]
is a supporting-side condition. Thus
$t=(\sqrt5-1)/2$. In each of the three cases the edge square $L$
fixes the circumradius square $L/(2-t)$, and $R$ fixes every remaining
dot product. This proves uniqueness of the face distance matrix up
to cyclic order and reflection. In particular for a pentagon the
diagonal-square to side-square ratio is $2+t=(1+\sqrt5)^2/4$.

\begin{corollary}[One shape per row of the table]\label{cor:unique}
For each of the five pairs in Section~\ref{sec:numerical}, at most
one similarity class of Platonic polyhedra exists.
\end{corollary}
\begin{proof}
Section~\ref{sec:maps} gives a face-incidence isomorphism between
any two members of the same case. Rescale one common edge length
to match. The preceding calculation shows that a convex regular
$p$-gon, for $p=3,4,5$, is determined by its side length and number
of sides. Corresponding facets are therefore congruent, and
Theorem~\ref{thm:cauchy} applies.
\end{proof}

The arm induction above is finite in general. For this application
its original links have at most five sides, and the midpoint cuts
introduce at most two more vertices. It can therefore be unrolled
into finitely many small cases if that is preferable in a proof
checker. The important point is that rigidity has been proved, not
silently assumed from the five numerical patterns.
